% \begin{document}
\section*{Lecture \#15: Type Inference and Let-Polymorphism}

\section{Polymorphism in OCaml}

In System F, programmers write type abstractions, type applications, and type
annotations explicitly; e.g.\ using the polymorphic $\mathit{twice}$ at $\kw{int}$
requires $\mathit{twice}\,[\kw{int}]\ (\lambda n : \kw{int}.\, n + 1)\ 7$. Since
reconstructing such annotations is undecidable (Wells 1994), a compiler cannot infer
them in general. Yet in OCaml no types need be written and we still get
polymorphism:
\begin{verbatim}
let twice f x = f (f x)
\end{verbatim}
is inferred to have type \texttt{('a -> 'a) -> 'a -> 'a}, which corresponds to the
System F type $\forall X.\, (X \to X) \to X \to X$. Then
\texttt{twice (fun n -> n + 1) 7} and \texttt{twice not true} both work, with the
type variable instantiated automatically at $\kw{int}$ or $\kw{bool}$.

OCaml can do this because its polymorphism is more restricted than System F's. We
first see how to infer types for the simply-typed $\lambda$-calculus via constraint
generation and \emph{unification}, then extend this to \emph{let-polymorphism}, the
restricted polymorphism of OCaml and the ML family.

\section{Type Inference}

In the simply-typed $\lambda$-calculus each argument has an annotation
$\lambda x : \tau.\, e$. Without annotations, can we still type check by
\emph{inferring} (or \emph{reconstructing}) the missing types? Yes. Consider
\[
\lambda a.\, \lambda b.\, \lambda c.\ \ttkw{if}\ a\,(b + 1)\ \ttkw{then}\ b\ \ttkw{else}\ c.
\]
$b$ is used in an addition, so $b : \kw{int}$. $a$ is applied to $b + 1$, so $a$ is a
function whose argument type is $\kw{int}$. The result of $a\,(b+1)$ is a guard, so
$a$ returns $\kw{bool}$. Both branches have the same type, so $c$ has the type of
$b$, namely $\kw{int}$:
\[
\lambda a : \kw{int} \to \kw{bool}.\ \lambda b : \kw{int}.\ \lambda c : \kw{int}.\ \ttkw{if}\ a\,(b + 1)\ \ttkw{then}\ b\ \ttkw{else}\ c.
\]

\subsection{Constraint-based typing}

Given $\Gamma$ and $e$, the algorithm produces a type together with a set of
\emph{constraints} (equations between types) that must hold for $e$ to be
well-typed. We extend types with type variables $X, Y, \ldots$ (placeholders for
unknown types, \emph{not} bound by any $\forall$, unlike System F's) and allow
abstractions with or without annotations:
\begin{align*}
e &::= x \mid \lambda x : \tau.\, e \mid \lambda x.\, e \mid e_1\, e_2 \mid () \mid n \mid e_1 \oplus e_2
      \mid \ttkw{true} \mid \ttkw{false} \mid \ttkw{if}\ e_1\ \ttkw{then}\ e_2\ \ttkw{else}\ e_3 \\
\tau &::= \kw{unit} \mid \kw{int} \mid \kw{bool} \mid \tau_1 \to \tau_2 \mid X
\end{align*}
The judgment $\Gamma \vdash e : \tau \mid C$ means $e$ has type $\tau$ in $\Gamma$
provided every constraint in $C$ is satisfied. Read bottom to top, the rules compute
$\tau$ and $C$ from $\Gamma$ and $e$.
\[
\inferrule*[Right=\textsc{CT-Var}]
  {\Gamma(x) = \tau}
  {\Gamma \vdash x : \tau \mid \emptyset}
\qquad
\inferrule*[Right=\textsc{CT-Int}]
  { }
  {\Gamma \vdash n : \kw{int} \mid \emptyset}
\qquad
\inferrule*[Right=\textsc{CT-True}]
  { }
  {\Gamma \vdash \ttkw{true} : \kw{bool} \mid \emptyset}
\]
\[
\inferrule*[Right=\textsc{CT-Arith}]
  {\Gamma \vdash e_1 : \tau_1 \mid C_1 \\ \Gamma \vdash e_2 : \tau_2 \mid C_2}
  {\Gamma \vdash e_1 \oplus e_2 : \kw{int} \mid C_1 \cup C_2 \cup \{\tau_1 = \kw{int},\, \tau_2 = \kw{int}\}}
\]
\[
\inferrule*[Right=\textsc{CT-If}]
  {\Gamma \vdash e_1 : \tau_1 \mid C_1 \\ \Gamma \vdash e_2 : \tau_2 \mid C_2 \\ \Gamma \vdash e_3 : \tau_3 \mid C_3}
  {\Gamma \vdash \ttkw{if}\ e_1\ \ttkw{then}\ e_2\ \ttkw{else}\ e_3 : \tau_2 \mid C_1 \cup C_2 \cup C_3 \cup \{\tau_1 = \kw{bool},\, \tau_2 = \tau_3\}}
\]
\[
\inferrule*[Right=\textsc{CT-AbsAnn}]
  {\Gamma, x : \tau_1 \vdash e : \tau_2 \mid C}
  {\Gamma \vdash \lambda x : \tau_1.\, e : \tau_1 \to \tau_2 \mid C}
\qquad
\inferrule*[Right=\textsc{CT-Abs}]
  {\Gamma, x : X \vdash e : \tau \mid C \\ X \text{ fresh}}
  {\Gamma \vdash \lambda x.\, e : X \to \tau \mid C}
\]
\[
\inferrule*[Right=\textsc{CT-App}]
  {\Gamma \vdash e_1 : \tau_1 \mid C_1 \\ \Gamma \vdash e_2 : \tau_2 \mid C_2 \\ X \text{ fresh}}
  {\Gamma \vdash e_1\, e_2 : X \mid C_1 \cup C_2 \cup \{\tau_1 = \tau_2 \to X\}}
\]
The rules for $()$ and $\ttkw{false}$ are like \textsc{CT-True}. In
\textsc{CT-AbsAnn} the annotation gives the argument type; in \textsc{CT-Abs} we
invent a fresh $X$; in \textsc{CT-App} we invent a fresh $X$ for the result and
require $e_1$ to be a function from the type of $e_2$ to $X$. ``Fresh'' means $X$
appears nowhere else in the derivation.

For our example, \textsc{CT-Abs} gives $a, b, c$ fresh types $X_a, X_b, X_c$.
\textsc{CT-Arith} on $b + 1$ gives $\kw{int}$ with constraints
$\{X_b = \kw{int},\, \kw{int} = \kw{int}\}$. \textsc{CT-App} on $a\,(b+1)$ gives a
fresh $Y$ and the constraint $X_a = \kw{int} \to Y$. \textsc{CT-If} gives type $X_b$
and adds $Y = \kw{bool}$ and $X_b = X_c$. So the whole expression has type
\[
X_a \to X_b \to X_c \to X_b
\]
subject to
\[
C = \{X_b = \kw{int},\ \kw{int} = \kw{int},\ X_a = \kw{int} \to Y,\ Y = \kw{bool},\ X_b = X_c\}.
\]

\subsection{Unification}

A \emph{type substitution} $\sigma$ is a finite map from type variables to types,
e.g.\ $[X \mapsto \kw{int},\, Y \mapsto \kw{int} \to \kw{int}]$. Applying it replaces
type variables simultaneously:
\[
\sigma(X) =
\begin{cases}
\tau & \text{if } X \mapsto \tau \in \sigma \\
X & \text{if } X \text{ is not in the domain of } \sigma
\end{cases}
\qquad
\sigma(\tau_1 \to \tau_2) = \sigma(\tau_1) \to \sigma(\tau_2),
\]
and $\sigma(\tau) = \tau$ for $\kw{unit}, \kw{int}, \kw{bool}$. As nothing binds
these type variables, capture is not a concern. Composition is
$(\sigma_1 \circ \sigma_2)(\tau) = \sigma_1(\sigma_2(\tau))$.

$\sigma$ \emph{unifies} a constraint $\tau_1 = \tau_2$ if $\sigma(\tau_1) =
\sigma(\tau_2)$, and \emph{satisfies} (unifies) a set $C$ if it unifies every
constraint in $C$. E.g.\ $[X \mapsto \kw{int},\, Y \mapsto \kw{int} \to \kw{int}]$
unifies $X \to (X \to \kw{int}) = \kw{int} \to Y$, as both sides become
$\kw{int} \to (\kw{int} \to \kw{int})$.

If $\Gamma \vdash e : \tau \mid C$ and $\sigma$ satisfies $C$, then $e$ has type
$\sigma(\tau)$ in $\sigma(\Gamma)$; if no substitution satisfies $C$, $e$ is not
typable. So we need a substitution satisfying $C$:
\begin{align*}
\unify(\emptyset) &= [\,] \quad \text{(the empty substitution)} \\
\unify(\{\tau_1 = \tau_2\} \cup C) &=
\begin{cases}
\unify(C) & \text{if } \tau_1 = \tau_2 \\
\unify(\tsub{C}{\tau_2}{X}) \circ [X \mapsto \tau_2] & \text{if } \tau_1 = X \text{ and } X \notin \ftv(\tau_2) \\
\unify(\tsub{C}{\tau_1}{X}) \circ [X \mapsto \tau_1] & \text{if } \tau_2 = X \text{ and } X \notin \ftv(\tau_1) \\
\unify(C \cup \{\tau_1' = \tau_2',\, \tau_1'' = \tau_2''\}) & \text{if } \tau_1 = \tau_1' \to \tau_1'' \text{ and } \tau_2 = \tau_2' \to \tau_2'' \\
\text{fail} & \text{otherwise}
\end{cases}
\end{align*}
Here $\ftv(\tau)$ is the set of type variables occurring in $\tau$. The algorithm
fails when asked to unify two different base types, or a base type with a function
type.

On our example: $X_b = \kw{int}$ yields $[X_b \mapsto \kw{int}]$ and turns
$X_b = X_c$ into $\kw{int} = X_c$; $\kw{int} = \kw{int}$ is already satisfied;
$X_a = \kw{int} \to Y$ yields $[X_a \mapsto \kw{int} \to Y]$; $Y = \kw{bool}$ yields
$[Y \mapsto \kw{bool}]$; finally $\kw{int} = X_c$ yields $[X_c \mapsto \kw{int}]$.
Applying the composition to $X_a \to X_b \to X_c \to X_b$ gives
\[
(\kw{int} \to \kw{bool}) \to \kw{int} \to \kw{int} \to \kw{int},
\]
matching the informal result. The algorithm is due to Robinson (1965), applied to type
inference by Hindley (1969) and Milner (1978).

\thm{Unification}{
The algorithm $\unify(C)$ always terminates. It succeeds if and only if $C$ has a
solution, and in that case $\unify(C)$ is a most general solution: for every solution
$\sigma'$ of $C$ there is a substitution $\sigma''$ such that
$\sigma' = \sigma'' \circ \unify(C)$.
}

Thus the inferred type is the most general one: every other type of the expression is
an instance of it. For $\lambda x.\, x$ we get $X \to X$ with no constraints, so
$\unify$ returns the empty substitution; every other type, e.g.\ $\kw{int} \to
\kw{int}$, is an instance of $X \to X$.

\subsection{The occurs check}

The check $X \notin \ftv(\tau)$ before mapping $X$ to $\tau$ is the \emph{occurs
check}. It prevents cyclic substitutions such as $[X \mapsto X \to Y]$, which make
sense only with infinite types. For $\lambda x.\, x\,x$, \textsc{CT-Abs} gives $x$ the
type $X$ and \textsc{CT-App} generates $X = X \to Y$; the occurs check makes
$\unify$ fail, so $\lambda x.\, x\,x$ is not typable (as in the simply-typed
$\lambda$-calculus). This is the source of OCaml's error ``The type variable 'a
occurs inside 'a -> 'b'' for \texttt{fun x -> x x}.

\section{Let-Polymorphism}

Unification finds the most general type, which may contain unconstrained type
variables; e.g.\ $\lambda x.\, x : X \to X$. We would like to use such a function at
several types:
\[
\ttkw{let}\ \mathit{id} = \lambda x.\, x\ \ttkw{in}\ \ttkw{if}\ \mathit{id}\ \ttkw{true}\ \ttkw{then}\ \mathit{id}\ 1\ \ttkw{else}\ 0
\]
(where $\ttkw{let}\ x = e_1\ \ttkw{in}\ e_2$ abbreviates $(\lambda x.\, e_2)\, e_1$).
With the rules so far, $\mathit{id}$ has the single type $X \to X$, and its two uses
generate $X = \kw{bool}$ and $X = \kw{int}$, which are unsatisfiable together.

The idea: at a $\ttkw{let}$ we know the whole definition, so we can generalize its
type to $\forall X.\, X \to X$ and instantiate it with fresh type variables at each
use. To keep inference decidable, $\forall$ is only allowed at the outside. A
\emph{type scheme} is
\[
s ::= \tau \mid \forall X.\, s,
\]
and contexts map variables to type schemes. The types $\tau$ contain no $\forall$ and
are called \emph{monotypes}.

\subsection{Typing rules}

For variables, abstractions, and $\ttkw{let}$ (other rules as in the simply-typed
$\lambda$-calculus):
\[
\inferrule*[Right=\textsc{T-Var}]
  {\Gamma(x) = \forall X_1 \ldots X_n.\, \tau}
  {\Gamma \vdash x : \tsub{\tau}{\tau_1/X_1, \ldots, \tau_n/X_n}{}}
\qquad
\inferrule*[Right=\textsc{T-Abs}]
  {\Gamma, x : \tau_1 \vdash e : \tau_2}
  {\Gamma \vdash \lambda x.\, e : \tau_1 \to \tau_2}
\]
\[
\inferrule*[Right=\textsc{T-Let}]
  {\Gamma \vdash e_1 : \tau_1 \\ \Gamma, x : \gen(\Gamma, \tau_1) \vdash e_2 : \tau_2}
  {\Gamma \vdash \ttkw{let}\ x = e_1\ \ttkw{in}\ e_2 : \tau_2}
\]
where $\gen(\Gamma, \tau) = \forall X_1 \ldots X_n.\, \tau$ and $X_1, \ldots, X_n$ are
the type variables free in $\tau$ but not free in $\Gamma$. We cannot generalize over
variables free in $\Gamma$: they stand for types of variables still in scope, which
are not yet known. \textsc{T-Var} instantiates the quantified variables with
arbitrary monotypes, and \textsc{T-Abs} gives the argument a monotype.

So $\mathit{id}$ gets $\forall X.\, X \to X$ and its two uses instantiate $X$ with
$\kw{bool}$ and $\kw{int}$; the whole expression has type $\kw{int}$.

$\ttkw{let}$ is no longer just an abbreviation:
\[
(\lambda \mathit{id}.\ \ttkw{if}\ \mathit{id}\ \ttkw{true}\ \ttkw{then}\ \mathit{id}\ 1\ \ttkw{else}\ 0)\ (\lambda x.\, x)
\]
is not typable, since the $\lambda$-bound $\mathit{id}$ must have a single monotype.
It is well-typed in System F, but only at $(\forall X.\, X \to X) \to \kw{int}$,
with a $\forall$ to the left of an arrow. Restricting $\forall$ to the outside of
type schemes is \emph{prenex polymorphism}; it makes inference decidable, and every
typable expression has a \emph{principal type scheme} of which all its other types are
instances.

\subsection{Inference}

Combine constraint generation with unification. To infer $\ttkw{let}\ x = e_1\
\ttkw{in}\ e_2$ in $\Gamma$:
\begin{enumerate}
  \item Generate constraints for $e_1$: $\Gamma \vdash e_1 : \tau_1 \mid C_1$.
  \item Solve them immediately: $\sigma = \unify(C_1)$.
  \item Generalize: give $x$ the type scheme $\gen(\sigma(\Gamma), \sigma(\tau_1))$.
  \item Infer $e_2$ in $\sigma(\Gamma)$ extended with this scheme for $x$.
\end{enumerate}
At each use of a variable with scheme $\forall X_1 \ldots X_n.\, \tau$, replace
$X_1, \ldots, X_n$ by fresh type variables. Constraints for $e_1$ must be solved
before generalizing, otherwise we would not know which variables are truly
unconstrained.

This is essentially Hindley--Milner inference (Algorithm~W), the core of type
inference in OCaml, Haskell, and others. It is fast in practice, though in the worst
case nested $\ttkw{let}$s can produce types of size exponential in the program.

\subsection{The value restriction}

With references, generalizing every $\ttkw{let}$ is unsound:
\begin{verbatim}
let r = ref (fun x -> x) in
r := (fun n -> n + 1);
(!r) true
\end{verbatim}
If $r$ were generalized to $\forall X.\, (X \to X)\ \kw{ref}$, the assignment could
use it at $\kw{int}$ and the dereference at $\kw{bool}$, applying an integer function
to a Boolean. OCaml's \emph{value restriction} generalizes
$\ttkw{let}\ x = e_1\ \ttkw{in}\ e_2$ only when $e_1$ is a syntactic value (such as a
function). Otherwise the type variables stay \emph{weak} and the first use fixes
them: above, the assignment fixes $r : (\kw{int} \to \kw{int})\ \kw{ref}$ and the
application to \texttt{true} is rejected.

\subsection{Weak type variables}

OCaml prints weak variables as \texttt{'\_weak1}, \texttt{'\_weak2}, \ldots: whereas
\texttt{'a} means ``for all types'', \texttt{'\_weak1} means one particular, not yet
determined type.
\begin{verbatim}
# let r = ref (fun x -> x);;
val r : ('_weak1 -> '_weak1) ref
# r := (fun n -> n + 1);;
- : unit = ()
# r;;
- : (int -> int) ref
# (!r) true;;
Error: This expression has type bool but an expression was expected of type int
\end{verbatim}
In terms of our algorithm, the value restriction solves the constraints for $e_1$ as
usual but skips $\gen$. A weak type variable is an unsolved type variable free in the
typing context; later constraints can solve it, but only once, since it stands for a
single type. This makes the restriction sound: a reference can only ever hold values
of one type.

The restriction also applies when nothing is mutable, since an application is not a
value:
\begin{verbatim}
# let ids = List.map (fun x -> x);;
val ids : '_weak1 list -> '_weak1 list
\end{verbatim}
The fix is to $\eta$-expand, making the definition a function (hence a value):
\begin{verbatim}
# let ids l = List.map (fun x -> x) l;;
val ids : 'a list -> 'a list
\end{verbatim}
A weak variable cannot escape a compilation unit: if nothing determines it, the
compiler reports that the type contains type variables that cannot be generalized.
OCaml in fact uses a \emph{relaxed value restriction} (Garrigue, 2004), which also
generalizes variables appearing only in covariant positions; e.g.\
\texttt{let x = List.rev []} has type \texttt{'a list}, since \texttt{'a} does not
appear under \texttt{ref} or to the left of an arrow.

\section{Summary}

System F gives fully general polymorphism but requires type annotations, as inference
is undecidable. For the simply-typed $\lambda$-calculus we infer types by generating
constraints and solving them by unification. Let-polymorphism extends this to a
restricted polymorphism with $\forall$ only at the outside of type schemes,
generalized at $\ttkw{let}$ and instantiated at each use. Inference is decidable and
finds principal types, which is why OCaml can infer polymorphic types for us.
% \end{document}
